% [ILLUSTRATION À GÉNÉRER : aquarelle d'une table autour de laquelle des salariés se partagent équitablement des pièces ou un gâteau — à créer via Midjourney puis insérer avec \insertImage]

\chapter{L'argent, le test ultime}

\lettrine[lines=2]{V}{ous} voulez savoir si une entreprise est vraiment libérée ? Ne regardez pas ses réunions, ses valeurs affichées ni son organigramme aboli. Posez une seule question : \emph{comment s'y décident les salaires ?}

L'argent est le test ultime du modèle, pour une raison simple : c'est le sujet où la confiance coûte le plus cher à accorder. Un dirigeant peut lâcher les horaires, les congés, le choix des outils, l'organisation des équipes — et garder la main sur les rémunérations. Tant qu'il la garde, le pouvoir n'a pas bougé d'un centimètre : il s'est juste fait discret. J'ai vécu les deux côtés de ce test, dans deux entreprises différentes. Voici ce que ça donne en vrai.

\section{La transparence des salaires : nécessaire, pas suffisante}

Dans la première entreprise libérée où j'ai travaillé, les salaires étaient totalement transparents — chacun pouvait connaître la rémunération de chacun. Sur le papier, c'est l'audace maximale, celle que la plupart des dirigeants classiques jugent impensable. Et je peux témoigner de son premier effet : la transparence assainit. Fini les rumeurs de couloir, les soupçons d'injustice fantasmée, les négociations secrètes qui récompensent les plus culottés plutôt que les plus utiles.

Mais j'ai aussi vu ses limites, et elles m'ont appris quelque chose d'essentiel. D'abord, cette transparence s'arrêtait aux frontières de notre site — le reste du groupe vivait à l'ancienne. Une transparence locale dans une organisation qui ne l'est pas, c'est une promesse à moitié tenue, et les demi-promesses fabriquent plus de frustration que l'opacité franche. Ensuite, et surtout : les \emph{augmentations}, elles, restaient décidées de manière classique, quelque part au-dessus de nous. Et il n'y avait pas de réelle répartition des richesses créées.

Voilà la leçon, et elle mérite d'être encadrée : \textbf{voir les salaires n'est pas un pouvoir ; c'est décider des salaires qui en est un.} La transparence sans processus de décision partagé, c'est une vitrine — on peut regarder, pas toucher. Elle est nécessaire, car aucune décision collective n'est possible sur des chiffres cachés. Elle n'est pas suffisante.

Des entreprises ont poussé la logique plus loin que nous. Buffer, l'éditeur américain de logiciels, publie depuis 2013 non seulement tous les salaires de ses équipes — dirigeant compris — mais la \emph{formule} qui les calcule, ouverte aux commentaires de tous\footnote{Gascoigne, J. (2013). Introducing Open Salaries at Buffer: Our Transparent Formula and All Individual Salaries. Buffer (buffer.com).}. La formule est discutable ; c'est justement le but : on ne peut discuter que ce qui est écrit.

\section{Qui décide des augmentations ?}

Dans ma propre société, nous avons donc attaqué le problème par l'autre bout : le processus de décision.

Concrètement, voici comment une augmentation se passait chez nous. Le salarié la demandait lors de son entretien individuel — jusque-là, rien d'exotique. Mais ensuite, la demande n'atterrissait pas sur le bureau du patron : elle était étudiée par un trio de trois salariés, qui l'examinait avec les vraies contraintes sur la table — l'état de la trésorerie, les prix du marché, la cohérence avec la grille transparente de l'entreprise. Le trio construisait une proposition, la présentait au salarié, et ce n'est qu'ensuite qu'elle m'arrivait — pour application, pas pour arbitrage.

Je veux insister sur ce dernier mot, parce que toute la philosophie du système est là : mon rôle était d'\emph{appliquer} la décision du collectif, pas de la valider. Et rappelez-vous le chapitre sur les leaders : c'est ce même processus, construit par les équipes pendant que je me retenais d'intervenir, qui m'a donné la plus belle preuve que la confiance fonctionne — des pairs qui connaissent la trésorerie ne décident pas plus mal que le patron. Ils décident mieux, parce qu'ils décident sous le regard de tous.

Nous n'avions rien inventé, d'ailleurs. Chez Morning Star, le transformateur de tomates californien, chaque salarié propose lui-même sa rémunération, examinée par un comité de pairs élus qui rend un avis argumenté\footnote{Laloux, F. (2014). Reinventing Organizations: A Guide to Creating Organizations Inspired by the Next Stage of Human Consciousness. Nelson Parker.}. Les mécanismes varient ; le principe est le même : sortir la décision salariale du tête-à-tête asymétrique entre un salarié et son chef, pour la confier à un processus visible, outillé par les chiffres réels de l'entreprise.

Soyons honnêtes sur les conditions de réussite, parce qu'il y en a. Ce système exige que les équipes connaissent la réalité économique — un trio qui ignore la trésorerie décidera dans le vide ; c'est toute la formation « à la vie d'une entreprise » dont ce livre parle depuis le début. Il exige du temps : trois personnes qui instruisent sérieusement une demande, c'est un investissement. Et il expose ceux qui décident : dire non à un collègue est plus difficile que de laisser un chef le faire. Aucun de ces coûts n'est rédhibitoire. Mais celui qui les cache vous prépare une désillusion.

\section{Et la répartition des richesses ?}

Il reste l'étage au-dessus, celui dont on parle le moins : au-delà des salaires, à qui vont les fruits ?

C'est la question qui fâche, parce qu'elle déborde du management vers la propriété. Une entreprise peut avoir les salaires les plus transparents et les processus les plus collectifs du monde — si les bénéfices et le capital suivent les circuits classiques, la libération s'arrête à la ligne « résultat » du compte d'exploitation. Dans ma première expérience, c'était exactement le cas : pas de réelle répartition des richesses créées. On peut appeler ça une entreprise libérée ; les salariés, eux, font le calcul.

Des réponses existent, à des degrés divers d'audace. L'intéressement et la participation, outils bien français, sont le premier étage — insuffisant seul, mais réel. Ricardo Semler, chez Semco, avait fait du partage des profits et de la transparence comptable les piliers de sa démocratie industrielle dès les années 1980\footnote{Semler, R. (1993). Maverick: The Success Story Behind the World's Most Unusual Workplace. Warner Books.}. Et tout au bout du chemin, il y a la réponse statutaire — celle des coopératives, où les salariés possèdent le capital. On y reviendra dans la conclusion, parce que c'était mon horizon personnel, et que je n'ai pas eu le temps de l'atteindre.

En attendant, si vous ne devez retenir qu'une chose de ce chapitre, que ce soit celle-ci : dans une entreprise libérée, l'argent n'est pas un sujet à part. C'est le sujet-vérité. Tout ce que le modèle promet — confiance, transparence, pouvoir partagé — passe ou casse à l'endroit précis où l'on décide qui gagne quoi. Une libération qui contourne ce sujet n'est pas prudente.

Elle est inachevée.
